\justifying
\NSFCBodyText

\subsubsection{研究背景与意义}

大模型智能体通过调用工具、读取环境反馈和调整行动计划，将语言理解能力转化为多步任务的执行能力。ReAct 将推理与行动交替组织，Toolformer 探索从数据中学习工具调用，使模型能够借助外部计算与信息完成任务\cite{react,toolformer}。与单次问答相比，智能体的输出还包括会改变后续状态的动作。一个语句流畅的回答可能对应错误的工具参数，一次接口返回成功也不意味着用户目标已经实现。因此，智能体的能力需要依据行动过程及最终状态评价，不能只看回答质量。

实际任务环境会持续变化：同类请求的组成改变，新工具加入，旧接口的字段或规则调整，反馈又可能延迟、缺失或出错。这些变化不仅降低当前任务成功率，还会改变适配数据的质量和模型置信度的含义。AgentBench、ToolLLM 和 $\tau$-bench 从多环境任务、工具调用及用户交互等角度推动了智能体评测\cite{agentbench,toolllm,taubench}，但在连续任务中，学会新工具、保留旧能力和避免错误执行需要同时面对资源约束。每次变化都重新训练全部参数，难以支持频繁适应；持续增加独立适配器，则会把训练效率问题转移为存储、路由与维护问题。

本项目题为\textbf{面向动态任务的大模型智能体参数高效自适应后训练与可信校准方法研究}，英文名称为 Parameter-Efficient Adaptive and Calibrated Training for LLM Agents under Dynamic Tasks，简称 PACT-Agent，面向 CCF--百度松果基金所关注的大模型后训练与 Agent 技术方向。项目以文本工具智能体为研究对象，考察固定资源预算内的持续适应机制。公开能源时序任务提供非平稳性和反馈延迟的补充验证环境，主体评价仍围绕通用工具任务展开。研究价值在于弄清何种任务变化值得增加参数更新、何种轨迹值得用于学习，以及策略更新后原有的行动置信度在何种条件下失效，从而为资源受限的持续后训练提供可验证的依据。

\subsubsection{国内外研究现状及发展动态分析}

\textbf{参数高效学习已具备成熟起点，跨任务预算分配仍需结合能力变化考察。}LoRA 冻结预训练主干，以低秩增量完成任务适配\cite{lora}；QLoRA 结合量化主干与低秩训练，降低显存占用\cite{qlora}。AdaLoRA 根据参数重要性在权重矩阵之间重新分配预算，说明不同层和模块的更新价值并不相同\cite{adalora}；DoRA 从权重幅度与方向分解改进低秩适配的表达能力\cite{dora}。这些研究表明，固定低秩配置并非唯一选择，也意味着“动态调整秩”不能直接作为新的科学贡献。对连续任务而言，需要进一步区分当前任务的适配收益、历史能力干扰，以及为估计这些收益所额外付出的计算成本。只有在累计存储与总计算预算相同的条件下，才能判断任务变化信息是否带来了有效的参数分配依据。

\textbf{后训练正在由单轮指令扩展到交互轨迹，监督质量成为关键变量。}InstructGPT 研究示范与偏好反馈对模型行为的影响\cite{instructgpt}，Self-Instruct 利用模型生成并筛选指令样本\cite{selfinstruct}，LIMA 则显示高质量样本值得单独考察\cite{lima}。DPO 将偏好学习转化为直接的策略优化目标，为轻量后训练提供了方法基础\cite{dpo}。在智能体场景中，AgentTuning 通过交互轨迹与通用指令混合训练增强智能体能力\cite{agenttuning}；Agent Lightning 将执行与训练解耦，以轨迹接口连接复杂交互和学习过程\cite{lightning}。AgentFlow 进一步在多轮交互中优化规划器，并结合执行、验证和记忆模块\cite{agentflow}。2026年的 Agent Lightning v1.0 还讨论了运行控制器参与训练时的数据拼接与优化细节\cite{lightningv1}。因此，建立多模块闭环本身已有充分先例，本项目需将研究重点放在证据对学习决策的作用及其可识别边界。

轨迹质量并不等价于最终成功标签。错误中间步骤可能被后续动作偶然修正，工具异常也可能被自然语言反馈掩盖。数学推理中的验证器研究说明，结果验证可以用于选择候选解答\cite{verifiers}；过程监督研究进一步比较了中间步骤与最终结果监督的差异\cite{stepverify}。但其可验证性依赖任务条件，不能将数学验证器的效果直接移用于任意接口。模型不确定性同样具有两面性：低置信轨迹可能包含重要的困难样本，也可能是错误反馈；高置信轨迹可能正确，也可能只是生成器与验证器共同偏差的结果。如何区分“执行正确”与“用于训练后有益”，是自动轨迹学习需要进一步回答的问题。

\textbf{持续学习提供了抗遗忘方法，但参数记忆与行为可靠性的关系尚需检验。}EWC 通过保护历史重要参数减轻遗忘\cite{ewc}；GEM 以情景记忆约束梯度并刻画前后向迁移\cite{gem}；LAMOL 通过生成旧任务样本支持语言持续学习\cite{lamol}。面向语言模型，O-LoRA 研究正交低秩子空间，LoRAMoE 通过低秩专家与路由缓解知识遗忘\cite{olora,loramoe}。Reflexion 则在不更新模型参数的情况下，以文字反思与情景记忆改善后续尝试\cite{reflexion}。这些工作为本项目提供直接对照，也提示必须区分参数更新、上下文记忆和额外尝试次数带来的收益。对工具智能体，旧答案仍然正确并不保证旧动作仍然符合更新后的接口规则；历史轨迹是否值得重放，需要同时检查能力覆盖和工具版本有效性。

\textbf{置信校准与选择性预测能够刻画可靠性，行动场景对信息时点提出更严格要求。}现代神经网络可能存在置信度与正确率不匹配的问题，温度缩放提供了简洁的后处理基线\cite{calibration}。Semantic Uncertainty 从语义等价角度刻画自然语言生成的不确定性\cite{semantic}。然而，多个表达相近的回答可能对应不同工具参数或不同环境后果。行动校准需要明确预测对象：模型估计的是候选动作满足前置条件的概率，还是整个任务最终成功的概率；两者不能混用。校准特征还必须在动作执行前可取得，不能将执行后才知道的验证结果回填为决策依据。

选择性预测通过拒绝部分样本改变风险与覆盖率的关系\cite{selective}。共形风险控制研究在明确条件下约束特定损失的期望，适应性共形推断研究分布变化中的长期覆盖\cite{conformalrisk,adaptiveconformal}。这些结论各有统计对象和假设，不意味着任意工具变化下的每个动作都能获得安全保证。WebArena 等可复现交互环境为功能正确性验证提供基础\cite{webarena}，但真正评价持续校准，还需要将策略版本、校准数据时点和被拒绝动作的潜在结果分开记录。否则，系统可能仅因拒绝更多任务而表现出较低错误率。

\subsubsection{现有研究的不足}

综上，现有研究分别推进了参数效率、轨迹监督、抗遗忘和置信校准。本项目关注这些方法进入动态任务序列后出现的三个具体问题。首先，任务变化强度并不直接等于参数需求，反馈噪声也可能被误判为能力不足；如果资源分配只依据当前损失，额外更新可能强化错误轨迹并干扰旧能力。其次，轨迹筛选与能力保持存在联系，过度保留容易成功的轨迹会缩小能力覆盖，简单重放旧轨迹又可能传播已失效的工具规则。再次，后训练改变了策略及其访问状态，旧校准器的误差不一定保持稳定；只有已执行动作留下结果时，还会产生选择偏差。

这三个问题不能仅靠增加模块数量解决。需要在累计参数、轨迹记忆和计算成本均受约束的条件下，区分任务漂移、验证证据与真实能力变化的作用，并观察这些作用如何传递到行动风险。项目因此以任务变化、轨迹证据和校准时点为共同研究线索，使背景问题能够对应到具体方法和验证终点。

\subsubsection{关键科学问题与研究切入点}

本项目的中心科学问题是：\textbf{在固定总资源预算和有限记忆下，轨迹证据如何影响参数适配、历史能力保持与行动风险之间的关系，其有效条件和失效边界是什么？}进一步提出三个子问题：任务状态与候选参数更新的边际收益之间是否存在可利用的稳定关系？执行证据、验证器反馈与不确定性如何共同区分轨迹正确性和训练价值？策略及工具版本改变后，历史记忆与独立校准如何影响行动选择的可靠性？

由此形成待检验的中心假设：在反馈可审计、任务变化具有一定可识别性且资源预算匹配的条件下，利用证据协调更新、记忆和校准，能够在预先规定的覆盖率、遗忘和行动风险容限内提高任务成功率；这一收益会随证据失真、反馈延迟和任务差异增大而减弱。若收益在匹配生成、验证和训练成本后消失，或只能通过持续扩充适配器与降低覆盖率获得，则该假设不成立。

研究将保留四条相互衔接的路线：任务感知的参数高效适配、证据驱动的可信轨迹学习、有界记忆下的持续抗遗忘、独立校准下的风险感知执行。各路线首先对其直接效果进行检验，再通过同一任务序列上的组合对照研究相互作用。预期贡献包括可比较的持续后训练方法、对证据失效情形的解释，以及包含成本、遗忘和风险覆盖率的可复现评价方案；不预设所有指标均同时改善，也不将尚未开展的实验写为已有结论。
